\section{Related works}
\label{sec:related_works}

\subsection*{3D Open Vocabulary Scene Understanding} Recent works~\cite{zhang2023clip, chen2023clip2scene, peng2023openscene3dsceneunderstanding, tschernezki22neural, goel2023interactive} have explored 2D vision-language models such as CLIP~\cite{Radford2021LearningTV}, OpenSeg~\cite{ghiasi2022scaling}, and OVSeg~\cite{liang2023open}, and foundation models such as SAM~\cite{kirillov2023segment} to ground 2D language-aligned features in 3D space. Initial approaches~\cite{kerr2023lerf, kobayashi2022decomposing, Haque_2023_ICCV} demonstrated the feasibility of distilling language features into NeRF-based representations, though at the cost of slow training and rendering. More recent Gaussian-based methods have improved efficiency while following a common framework: extracting language-aligned feature maps from 2D views and grounding them to 3D space through multi-view optimization. LangSplat~\cite{qin2024langsplat} employs SAM to obtain masks at different granularity, then extracts CLIP image features for these regions. LEGaussians~\cite{shi2024language} combines DINO~\cite{caron2021emerging} and CLIP features to obtain hybrid dense language feature maps and FMGS~\cite{zuo2024fmgs} uses a similar combination to improve robustness. GOI~\cite{qu2024goi3dgaussiansoptimizable} leverages APE~\cite{shen2024aligning} to extract pixel-aligned features with fine boundaries.

However, incorporating the high-dimensional language features into 3D space presents unique challenges. It not only increases memory requirements but also significantly extends training time. Different approaches have been proposed to address this challenge: LangSplat~\cite{qin2024langsplat} employs a scene-specific autoencoder for feature compression, LEGaussians~\cite{shi2024language} uses quantization of semantic features, and Feature-3DGS~\cite{zhou2024feature} employs feature compression with a lightweight decoder. FMGS~\cite{zuo2024fmgs} introduces multi-resolution encoding for efficient semantic embedding and GOI~\cite{qu2024goi3dgaussiansoptimizable} utilizes a trainable feature clustering codebook.

In contrast to these methods, our approach supports arbitrary-dimensional uncompressed features without compromising training speed. Our method preserves semantic meaning without requiring additional preprocessing steps or encoder-decoder structures. This ensures open-vocabulary generalization to the same extent as the underlying feature extractor and enables a range of downstream applications using dynamic scenes or scene editing.

\subsection*{Uplifting 2D Features to 3D} 
While current methods~\cite{qin2024langsplat,shi2024language,peng20243d, kobayashi2022decomposing} focus on optimization-based feature lifting, several approaches also explore direct back-projection of 2D features to 3D space using pixel rays and depth information.
OpenGaussian~\cite{wu2024opengaussian} proposes an instance-level 3D-2D feature association method that links 3D points to 2D masks. Semantic Gaussians~\cite{guo2024semantic} utilizes depth to find corresponding 3D Gaussians along each pixel's ray. PLGS~\cite{wang2024plgsrobustpanopticlifting} initializes semantic anchor points by back-projecting pixels into 3D space using camera parameters. While not focusing on semantic understanding, FlashSplat~\cite{flashsplat} demonstrates that 2D mask rendering is essentially a linear function of Gaussian labels, enabling closed-form optimization through linear programming. 

In contrast, our method accurately models Gaussian Splatting forward rendering to lift 2D language features to 3D, while avoiding back-projection limitations such as occlusion and incomplete depth information. Table~\ref{tab:method-comparison} summarizes the key differences between our method and existing approaches. Unlike previous approaches, our method is training-free and does not require costly semantic feature rendering. While maintaining comparable performance with existing methods, it efficiently completes the lifting of 512-dimensional features to 3D in less than one minute.